\section{Introduction}\label{sec:intro}

Optimization models in process systems engineering, energy networks and
optimal control are often built from local pieces. A flowsheet links unit
models along streams, a multiperiod model links consecutive time stages, and a
network model links quantities at adjacent nodes. The objective is then a sum
of terms that each depend on a few variables, and the graph recording which
variables share a term often has small treewidth. For variables with finitely
many values, small treewidth is exactly what dynamic programming needs:
nonserial dynamic programming, bucket elimination and junction-tree message
passing find a global optimum in time exponential only in the width of a tree
decomposition \cite{BerteleBrioschi1972,Dechter1999,KollerFriedman2009}.

Mixed-integer nonlinear models have continuous variables and nonconvex
terms, and there the situation is less clear. Spatial branch-and-bound
solvers relax individual nonconvex terms by convex underestimators and branch
on variable domains \cite{TawarmalaniSahinidis2005,BelottiEtAl2009,
MisenerFloudas2014,BelottiEtAl2013}; they do not use a tree decomposition to
organize the search. Discretizing the continuous variables and running a
dynamic program gives neither a valid bound nor a running time that scales
well with the accuracy: a uniform grid of mesh $h$ on a unit interval has
$\Theta(1/h)$ points, so tables grow like $h^{-p}$ in the bag size $p$. Sparse
polynomial optimization exploits the same structure through linear and
semidefinite hierarchies \cite{WakiEtAl2006,Lasserre2006,BienstockMunoz2018};
the known bounds on their size are polynomial in the inverse accuracy
\cite{BienstockMunoz2018,KordaMagronRiosZertuche2025}.

Small treewidth alone does not make continuous nonconvex problems
tractable. Del Pia and Khajavirad \cite{DelPiaKhajavirad2026} recently showed
that box-constrained nonconvex quadratic programming (box QP) is solvable in
strongly polynomial time on forests, but strongly NP-hard at treewidth two even with
small integer coefficients, and that quartic minimization is strongly NP-hard
on paths. This paper asks which additional quantitative information makes
tree decompositions useful for \emph{certified} global optimization of
mixed-integer problems, and what fails when it is absent.

\subsection{Results}\label{sec:intro-results}
We consider $\min\{F(x):x\in X\}$, with optimal value $\OPT$, where
$F=\sum_af_a(x_{S_a})$ is a sum of factors, $X$ is a box in which some coordinates are integer, and a
tree decomposition of the scopes $S_a$ with maximum bag size $p$ is given.
Two quantities govern the results: upper bounds $L_i$ on the curvature of $F$
along each coordinate (for a quadratic, $L_i=\max\{H_{ii},0\}$), and
quadratic growth of $F$ away from a minimizer $x^*$. \emph{Point growth}
$F(x)-\OPT\ge g\norm{x-x^*}^2$ on $X$ gives the condition number
$\kappa=\max\{1,L/g\}$ with $L=\max_iL_i$. \emph{Weighted growth}
$F(x)-\OPT\ge\gamma\sum_iL_i(x_i-x_i^*)^2$ on $X$ gives
$\bar\kappa=\max\{1,1/\gamma\}$. Point growth implies weighted growth with
$\bar\kappa\le\kappa$, and $\bar\kappa$ does not change when continuous
coordinates are rescaled (Definition~\ref{def:growth}).

Choose finite grids in the coordinates and subtract, at each node $v$ of
coordinate $i$, the unary correction $L_iw^2/8$, where $w$ is the length of the
longest grid interval ending at $v$. The observation behind our certificate is that this
node-separable correction reproduces the vertex lower bound of Bajaj and
Hasan \cite{BajajHasan2020}: the minimum of the corrected objective over the
grid equals the best vertex bound, taken with per-coordinate curvature
bounds, over all cells of the product grid (Proposition~\ref{prop:cellwise}).
It is therefore a lower bound on $\min_XF$. Because the corrections are
unary, this minimum and all coordinate min-marginals are computed by ordinary
dynamic programming on the given decomposition, without enumerating cells.
The min-marginals give interval
bounds for optimality-based bound tightening from one pair of message passes:
an interval whose endpoint min-marginals exceed the value of a known feasible
point contains no better point and is removed. The grids of the successive
stages, one bound and one feasible point form a \emph{path certificate},
which an independent program checks by recomputation
(Theorem~\ref{thm:certificate}). None of this uses convexity, uniqueness or
growth.

Growth enters only the running time. Under weighted growth, grids whose
spacing grows linearly away from the current corrected minimizer
(\emph{graded grids}), combined with this filtering, keep
$O(\theta^{-1}\log(n+2))$ nodes per coordinate at every stage, independently
of the accuracy, when the grading $\theta$ satisfies $8\bar\kappa\theta^2\le1$
(Lemma~\ref{lem:states}); for $\theta$ of order $1/\sqrt{\bar\kappa}$ this is
$O(\sqrt{\bar\kappa}\log(n+2))$. The algorithm CT
(Algorithm~\ref{alg:ct}) tries $\theta=2^{-\mu}$ for $\mu=2,3,\dots$ with a cap on
the grid size, so the growth constant need not be known. The following theorem
summarizes the main results; $I$ denotes the input length and $q\ge0$ an
integer.

\begin{theorem}[Main results]\label{thm:intro-main}
Let $F$ be an explicit rational quadratic on a mixed box $X$, given as a sum
of factors with a tree decomposition of maximum bag size $p$, and let
$L_i=\max\{0,\partial_{ii}F\}$.
\begin{enumerate}[label=(\alph*)]
\item On every instance, CT$(2^{-q})$ returns a feasible rational point $\hat
x$ and a path certificate showing $\beta\le\OPT\le F(\hat x)\le\beta+2^{-q}$.
The certificate is checked by recomputing one dynamic program per stage in
exact arithmetic, and its validity uses no convexity, uniqueness or growth
(Theorems~\ref{thm:certificate} and~\ref{thm:approx}).
\item If $F$ has weighted growth with condition number $\bar\kappa$, every grid
on which CT$(2^{-q})$ builds tables has $O(\sqrt{\bar\kappa}\log(n+2))$ nodes per
coordinate, independently of $q$, and CT$(2^{-q})$ uses at most
$f(p,\bar\kappa)(I+q+1)^5$ bit operations, where
$f(p,\kappa)=c_0(c_1p\sqrt\kappa)^p\kappa(1+\log_2\kappa)^2$ for absolute
constants $c_0,c_1$. The certificate has size, and is checked in time, within
the same bound. Neither the growth constant nor $\bar\kappa$ is an input
(Lemma~\ref{lem:states} and Theorem~\ref{thm:approx}).
\item The algorithm EX (Algorithm~\ref{alg:ex}) returns a global minimizer and
$\OPT$ on every instance, with a certificate whose validity uses neither
growth nor uniqueness. If $F$ has point growth with condition number
$\kappa$, EX uses at most $f_1(p,\kappa)(I+1)^{C_1}$ bit operations, for a
computable function $f_1$ and an absolute constant $C_1$
(Theorem~\ref{thm:exact}).
\end{enumerate}
\end{theorem}

To the best of our knowledge, no earlier result gives a running time of the
form $f(p,\kappa)\poly(I+q)$ for certified approximation, or
$f(p,\kappa)\poly(I)$ for exact solution, of nonconvex or mixed-integer box
quadratic programs (Section~\ref{sec:related}). Under point growth, the
approximation result extends to fixed-degree polynomial factors, with exact
output when all coordinates are integer (Corollary~\ref{cor:poly}); continuous polynomial
objectives can have irrational minimizers (Example~\ref{ex:polylimits}).
Exactness itself needs no uniqueness: explicit rational height bounds give a
snapping procedure that turns any feasible point with certified gap at most
$2^{-k}$, $k=\poly(I)+\max\{0,\log_2(1/g_S)\}$, into an exact minimizer,
where $g_S>0$ is a growth constant towards the optimal \emph{set}; such a
constant exists for every instance (Theorem~\ref{thm:transfer} and
Remark~\ref{rem:setgrowth}). The growth hypothesis of part~(c) serves only
to bound the time needed to reach this accuracy.

\paragraph{Parameterized form.}
The bound in part~(b) holds on every instance with weighted growth. The
number $\bar\kappa$ is a property of the instance, which the algorithm
neither receives nor computes. Since $f$ is increasing in $\kappa$, the bound
is at most $f(p,\lceil\bar\kappa\rceil)(I+q+1)^5$; in parameterized terms
this is a fixed-parameter bound in the parameters $p$ and
$\lceil\bar\kappa\rceil$, which the algorithm does not need to know
(Remark~\ref{rem:fpt}). The graded grids are what give this form: filtered
uniform grids can keep order $\sqrt{n\kappa}$ nodes per coordinate
(Corollary~\ref{cor:uniformgrid}), so their bag tables can have order
$(n\kappa)^{p/2}$ entries, which is polynomial for fixed $p$ but not of this
form.

\paragraph{Scope of the growth hypothesis.}
Growth restricts where nonconvexity can occur. For a box QP, point growth
forces the Hessian block of the continuous coordinates strictly inside their
bounds at the minimizer to be positive definite, and weighted growth bounds
this block below by $2\gamma\diag(L_i)$ (Lemma~\ref{lem:growthcert}). The
nonconvex instances covered by our bounds therefore have their negative
curvature in directions that involve active bounds or integer coordinates;
such instances can have exponentially many strict local minima
(Example~\ref{ex:family}). The condition number can be large: unless
P${}={}$NP, no polynomial in $I$ bounds $\kappa$ on the hard instances of
Theorem~\ref{thm:intro-lower}(a) below, on which $\log\kappa=O(I)$, because
CT would then solve them in polynomial time. We do not claim that $\kappa$ or
$\bar\kappa$ is moderate for the application models mentioned above. The certificates are
valid in every case; $\kappa$ and $\bar\kappa$ bound only the running time.

\paragraph{Exact messages.}
A chain with box widths $2^t-1$ (Example~\ref{ex:chain}) shows the difference
from exact parametric dynamic programming: its exact Bellman messages need at
least $2^{m-1}$ quadratic pieces (Proposition~\ref{lim:prop:messages}), while
CT certifies any accuracy $2^{-q}$ in time polynomial in $m$ and $q$. Our
implementation certifies a gap of $10^{-3}$ for $m=64$ with at most eleven
nodes per coordinate grid (Section~\ref{sec:computation}).

\paragraph{Lower bounds.}
The dependence on both parameters is necessary under standard complexity
assumptions: P${}\ne{}$NP; the exponential-time hypothesis (ETH) that, for
some $\delta>0$, 3-SAT with $m$ variables cannot be solved in time
$2^{\delta m}$ \cite{ImpagliazzoPaturiZane2001}; and its randomized version
(rETH), which makes the same claim for randomized algorithms.

\begin{theorem}[Lower bounds]\label{thm:intro-lower}
\leavevmode
\begin{enumerate}[label=(\alph*)]
\item Unless P${}={}$NP, no algorithm computes, for every rational continuous
box QP with a supplied tree decomposition of maximum bag size three and a
unique minimizer, $\OPT$ within $2^{-q}$ in time polynomial in
$I+q+\log\kappa$. The hard instances have $\log\kappa=O(I)$
(Proposition~\ref{lim:prop:unique} and Corollary~\ref{lim:cor:nopolylog}).
\item Computing $\OPT$ within $\frac12$ is NP-hard for box QPs with
$\bar\kappa\le\kappa\le2$, continuous or all-binary, by a reduction from
maximum independent set; under ETH it cannot be done in time
$2^{o(p)}I^{O(1)}$ (Proposition~\ref{prop:lbwidth}).
\item Under rETH there are no computable $f$ and computable nondecreasing
unbounded $\psi\ge1$ such that some algorithm outputs, for every integer box QP
with a supplied tree decomposition of maximum bag size $p$ and a unique
minimizer, a feasible point
within $\frac12$ of $\OPT$ in time $f(p)(\kappa I)^{p/\psi(p)}$
(Proposition~\ref{prop:lbproduct}).
\item For $p\ge1$ and $\kappa\ge8p$, every deterministic algorithm that
accesses $F$ only through its values and returns a point with gap below
$Lp/\kappa$ for all functions on $[0,1]^p$ with upper coordinate curvature $L$, a
unique minimizer and point growth $L/\kappa$ makes at least
$(\kappa/(8p))^{p/2}-1$ evaluations on one of them
(Proposition~\ref{lim:prop:oracle}).
\end{enumerate}
\end{theorem}

Thus, unless P${}={}$NP, neither parameter can be dropped and the dependence
on $\kappa$ cannot be polylogarithmic; under ETH the dependence on $p$ is
exponential even for $\kappa\le2$; and under rETH the exponent of $\kappa$
cannot be $o(p)$, already for integer box quadratics, whereas
Theorem~\ref{thm:intro-main}(b) has exponent $p/2+O(1)$. Because
$\bar\kappa\le\kappa$, the same statements hold for $\bar\kappa$. In the
value-oracle model of part~(d), CT attains the exponent $p/2$ up to a factor
$(C\sqrt p\log(p+2))^p$ and the number of stages
(Remark~\ref{lim:rem:oracle}). Section~\ref{sec:limits} also shows the
following: exact piecewise-quadratic messages can need exponentially many
pieces at bag size three and $\kappa\le80$; growth towards a continuum of
minimizers does not bound the grids produced by filtering; for smooth
functions with growth towards an unknown optimal set, an algorithm that
queries only values and derivatives needs a number of queries that grows
with the accuracy; a path
of local affine equalities makes the problem NP-hard at bag size three and
$\kappa=1$, by the running-sum encoding of Subset Sum of Bienstock and
Mu\~noz \cite{BienstockMunoz2018} (see also
\cite[Example~1.1]{CifuentesParrilo2016}); and matching separator
moments of any finite order does not give a local error bound in terms of
negative curvature. As far as we know, the expanding-box chain, the
unique-minimizer hardness with explicit growth (a refinement of
\cite[Remark~2]{DelPiaKhajavirad2026}), the rETH bound and the moment
obstruction are new; the other statements adapt standard constructions.

\paragraph{Conditional recourse.}
We call exact minimization over some variables, for fixed values of the
others, \emph{recourse}. The condition number uses diagonal curvature, so a
stiff convex coupling term inflates it even though the convex subproblem it
defines is easy. The corrected-grid bound and the growth analysis apply to
any value function $V(v)=\min_yF(v,y)$ whose inner feasible set does not
depend on $v$, because partial minimization preserves upper coordinate
curvatures and growth (Lemma~\ref{lem:valuefunction} and
Theorem~\ref{thm:valuefn}). Bag-local corrections are valid with exact or
certified recourse (Theorem~\ref{thm:cr-filter}), but not when the outside
coordinates are optimized on the grid: the required allowance then grows with
their number (Proposition~\ref{prop:star}). Private convex quadratic blocks
become value factors whose curvature is certified by piecewise-affine
responses, with no global overlay of the pieces (Theorem~\ref{thm:cv}); the
certified cancellation is the largest possible for the block
(Proposition~\ref{prop:vf-curv}). A residual problem that is concave in each
coordinate and submodular after sign changes is evaluated exactly, with a
flow certificate, by one minimum cut at each core point, so only the core is
discretized (Theorems~\ref{thm:cr-oracle} and~\ref{thm:cr-search}). For
balanced box quadratics the grid problem itself is a minimum cut, which
removes the width parameter: under point growth, certified approximation
takes $(n\kappa(I+q+1))^{O(1)}$ bit operations
(Theorem~\ref{thm:balanced}). Recourse does not reduce the parameter to the ratio of negative curvature to growth in
general (Proposition~\ref{prop:cv-limit}). Parametric quadratic programming,
cut representations of submodular quadratic functions and threshold
encodings of ordered labels are classical; what this paper adds is the
curvature and growth accounting that makes them valid inside the certified
grid algorithm.

\paragraph{Coupling constraints.}
Independent rounding, on which the certificate rests, breaks under coupling
constraints. Totally unimodular (TU) coupling constraints with mesh-aligned
data admit exactly feasible correlated rounding, which restores the certificate
and, for quadratics, exact output (Section~\ref{sec:constraints}). Under set
growth with finitely many optimal values per coordinate, the algorithm is
polynomial for fixed width when the condition number $\kappa_c$ (the ratio
of a curvature bound for the continuous variables to the set-growth
constant), the box widths measured in mesh units ($s/\eta$), and the number
$r$ of optimal values
per coordinate are polynomially bounded (Theorem~\ref{thm:tu-approx}). It is
not fixed-parameter tractable in $(p,\kappa_c)$, because its meshes are
uniform (Section~\ref{sec:tu-limits}), and the pseudopolynomial level-0 grid
cannot be removed unless P${}={}$NP (Proposition~\ref{lim:prop:constraints}).
Two natural non-uniform alternatives, misaligned coordinate grids and a
common graded pattern, give invalid bounds with curvature-only corrections
(Proposition~\ref{prop:tu-misaligned} and Example~\ref{ex:tu-sum}). Bienstock
and Mu\~noz approximate bounded-width polynomial problems with arbitrary
polynomial constraints by linear programs whose solutions are feasible within
a scaled tolerance \cite{BienstockMunoz2018}; our extension covers a narrower
class but is exactly feasible (Remark~\ref{rem:tu-bm}).

\paragraph{Several minimizers.}
With several minimizers, one graded grid can need exponential time even for
two isolated minimizers (Proposition~\ref{prop:twocenters}). When the optimal
set is finite, unions of uniform cells give grids with $O(r\sqrt{n\kappa_S})$
nodes per coordinate, where $r$ bounds the number of optimal values per
coordinate and $\kappa_S=\max\{1,L/g_S\}$ uses the growth constant $g_S$
towards the optimal set (Theorem~\ref{thm:cells}); for fixed $p$ the work is
polynomial in $n$, $r$, $\kappa_S$ and $I+q$. For coordinatewise concave
mixed-integer quadratics, the Bellman residuals of an endpoint dynamic program
describe the entire optimal set as a tree-structured constraint problem, from
which uniqueness, the dimension of the optimal set and, when it is finite,
the number of optimizers are computed in $2^{O(p)}\poly(I)$ bit operations
(Theorem~\ref{thm:endpointset} and Corollary~\ref{cor:facecsp}). This combines
Rosenberg's description of the optimal set of a multilinear function as a
union of faces \cite{Rosenberg1972} with zero-residual descriptions of all
optimal labelings of discrete problems
\cite{WainwrightJaakkolaWillsky2005,Werner2007}; what is new is the treatment
of strictly concave coordinates and mixed boxes, and the factored description
on the decomposition. Finally, for continuous box quadratics certified by a
diagonal Lagrangian multiplier, an exact optimizer and an exact description
of the optimal set by linear equations and complementarity conditions are
found in $f'(p,\kappa_S)\poly(I)$ bit operations, for a computable function
$f'$, without knowledge of $g_S$ (Theorem~\ref{thm:diagdiscovery}). The certificate is classical
\cite{JeyakumarRubinovWu2006,LiWuQuan2015}; what we use is that it is the
same at every optimizer.

Table~\ref{tab:results} lists the main results with their hypotheses and
bounds.

\begin{table}[tbp]
\centering
\small
\setlength{\tabcolsep}{3pt}
\caption{Main results. Bounds count bit operations unless stated otherwise;
$I$ is the input length and $2^{-q}$ the certified gap. Certificates are
valid without the growth hypotheses, which enter only the running times.
Notation: $\kappa(L',g')=\max\{1,L'/g'\}$; $f$ and $f_1$ are as in
Theorem~\ref{thm:intro-main}, $f'$ is a computable function, and $C,C_1$ are
absolute constants unless stated otherwise; in the TU row, $n_c$ is the number of continuous coordinates and $K_Z\le I$
bounds the number of values of a discrete coordinate.}\label{tab:results}
\begin{tabular}{@{}>{\raggedright\arraybackslash}p{2.2cm}
>{\raggedright\arraybackslash}p{4.2cm}
>{\raggedright\arraybackslash}p{1.8cm}
>{\raggedright\arraybackslash}p{4.0cm}
>{\raggedright\arraybackslash}p{3.4cm}@{}}
\toprule
Result & Setting and hypothesis & Parameters & Bound & Output\\
\midrule
Thm.~\ref{thm:approx} & mixed box QP with tree decomposition; weighted
growth & $p,\bar\kappa$ & $f(p,\bar\kappa)(I+q+1)^5$ & feasible point, path
certificate\\
Thm.~\ref{thm:exact} & as above; point growth & $p,\kappa$ &
$f_1(p,\kappa)(I+1)^{C_1}$ & exact minimizer and $\OPT$, certificate\\
Thm.~\ref{thm:valuefn} & value function $V$ whose value factors are
evaluated exactly in polynomial time; point growth of $V$ (exact output:
quadratic $F$ with point growth) &
$p$ of $V$, $\kappa_V=\kappa(L^V,g)$ &
$f(p,\kappa_V)\kappa_V^{C}(I+q+1)^{C}$; exact
$f(p,\kappa_V)\kappa_V^{C_1}(I+1)^{C_1}$; $C,C_1$ depend on the
polynomial time bounds &
feasible original point, certificate; exact optimizer\\
Thm.~\ref{thm:cv} & private convex QP blocks with leaf certificates;
growth of $V$ (exact output: of $F$) & $p,\kappa(L^+,g)$, $L^+$ reduced &
$f(p,\kappa)\kappa^{C}(I+q+1)^{C}$; exact
$f(p,\kappa)\kappa^{C_1}(I+1)^{C_1}$; $2^{O(p)}\poly(I)$ if $L^+=0$ &
feasible point, exact optimizer\\
Thms.~\ref{thm:cr-oracle}, \ref{thm:cr-search}, \ref{thm:cr-exact} &
residual concave in each coordinate and balanced; continuous core of $k$
coordinates; core growth & $k,\kappa_{\mathcal K}$ & one maximum flow per
query; exact $8^k(1+\sqrt{k\kappa_{\mathcal K}})^k\poly(I)$ & conditional
values with flow certificates; exact optimizer\\
Thm.~\ref{thm:balanced} & balanced box QP, no decomposition; point growth &
$n,\kappa$ & $(n\kappa(I+q+1))^{O(1)}$; exact $(n\kappa(I+1))^{O(1)}$ &
certificate by cuts and flows; exact minimizer\\
Thm.~\ref{thm:tu-approx} & TU coupling, mesh-aligned data; set growth, at
most $r$ optimal values per coordinate & $p,\kappa_c,s/\eta,r$ & $K^p$ table
entries per level, $O(I+q)$ levels, $K=\max\{K_Z,\,s/\eta+1,$
$r(5+2\sqrt{n_c\kappa_c})\}$ & exactly feasible point, certificate; exact
output for quadratics\\
Thm.~\ref{thm:cells} & finite optimal set; set growth & $p,r,\kappa_S$ &
$O(r\sqrt{n\kappa_S})$ nodes per coordinate; $K_S^p\poly(I+q)$; exact
$K_S^p\poly(I+\log\kappa_S)$ & certificate; exact minimizer\\
Thm.~\ref{thm:endpointset}, Cor.~\ref{cor:facecsp} & $H_{ii}\le0$ for all $i$,
mixed box; no growth & $p$ & $2^{O(p)}\poly(I)$ & all optimizers as a
tree-structured constraint problem\\
Thm.~\ref{thm:diagdiscovery} & continuous box QP with diagonal Lagrangian
certificate; set growth & $p,\kappa_S$ & $f'(p,\kappa_S)\poly(I)$ & optimizer
and description of the optimal set\\
\midrule
Cor.~\ref{lim:cor:nopolylog} & continuous box QP, $p=3$, unique minimizer & $\kappa$ & no
$\poly(I+q+\log\kappa)$ unless P${}={}$NP & lower bound\\
Prop.~\ref{prop:lbwidth} & box QP, $\kappa\le2$ & $p$ & no $2^{o(p)}I^{O(1)}$
under ETH & lower bound\\
Prop.~\ref{prop:lbproduct} & integer box QP, unique minimizer & $p,\kappa$ &
no $f'(p)(\kappa I)^{p/\psi(p)}$ with $\psi$ unbounded, under rETH & lower
bound\\
Prop.~\ref{lim:prop:oracle} & value oracle, one bag & $p,\kappa$ &
$(\kappa/(8p))^{p/2}-1$ evaluations & lower bound\\
\bottomrule
\end{tabular}
\end{table}

\paragraph{Organization.}
Section~\ref{sec:related} reviews related work. Section~\ref{sec:setting}
fixes the model, the curvature bounds and the growth conditions, and closes
with a table of notation. Section~\ref{sec:grids} introduces corrected grids,
filtering and path certificates. Section~\ref{sec:growth} proves the
accuracy-independent grid bound and part~(b) of
Theorem~\ref{thm:intro-main}, and Section~\ref{sec:exact} gives exact output,
part~(c). Sections~\ref{sec:recourse}, \ref{sec:constraints}
and~\ref{sec:optsets} treat conditional recourse, coupling constraints and
several minimizers. Section~\ref{sec:limits} proves
Theorem~\ref{thm:intro-lower} and the further limits,
Section~\ref{sec:computation} reports the implementation and experiments, and
Section~\ref{sec:conclusion} lists open problems. The appendices follow the
order of the sections. Besides longer proofs, they contain secondary results:
exact implicit output for polynomial objectives with boundary optima
(Appendix~\ref{app:boundary}), further results on convex and cut recourse,
among them exact output with a cut residual
(Appendices~\ref{app:recourse-convex} and~\ref{app:recourse-cuts}), the
diagonal certificate and the proximal stage of Section~\ref{sec:diagcert}
(Appendix~\ref{app:proximal}), and the moment obstruction
(Appendix~\ref{app:moments}).
